% `tex.web` §392-§397 and §476: matching a delimited argument, and the two
% things that make it more than a scan for a token run.
%
% §476: a `#{` at the end of a parameter text makes the brace the *last token
% of the delimiter*, not something after it. `\def\t12#101001#{…}` delimits
% `#1` by `01001{`, so the `01001` at the front of `0100101001001{` does not
% end the argument — no `{` follows it. We took the literal run alone and
% matched there, handing back an empty argument.
%
% §397: a failed partial match is not thrown away. A match may begin *inside*
% it, so the window is shifted one token at a time and re-tested — TeX walks
% `u` and `v` back over the matched tokens for exactly this. Discarding the
% partial and restarting at the current token never tries those starts, so even
% with the brace in the delimiter the `01001{` that begins four tokens later
% was invisible and the argument ran to the end of the file.
%
% Each fault alone gives a different wrong answer and together they gave a
% third. trip.tex line 161 is an `\ifdim` whose branch turns on the result.
\catcode`\{=1 \catcode`\}=2 \catcode`\#=6

\def\t12#101001#{-.#1pt}
\message{<brace-delim=\t120100101001001{\relax}>}

% The `#{` trick with no literals before it: read up to the brace, which stays.
\def\u#1#{[#1]}
\message{<plain-hash-brace=\u abc{\relax}>}

% A delimiter that overlaps itself, which is what needs the backtracking: the
% first `aab` starts inside the failed match of `aab` at the front.
\def\v#1aab{(#1)}
\message{<overlap=\v aaabX>}

% ...and an ordinary delimiter, unchanged.
\def\w#1XY{[#1]}
\message{<plain=\w abcXY>}
\end
